% Ad Atticum 15.27 --- headnote
% ad-atticum-15-27, 3 July 44 BC, in Arpinati

\textit{Cicero to Atticus from the Arpinum estate, 3 July 44 BC
--- Perseus dateline \textit{Scr. in Arpinati. v Non. Quint. a.
710 (44)}, written the day after 15.26 and a few hours before
15.28 of the same day. A short, intimate letter in three sections.
Cicero is glad Atticus has urged him to write to Sestius, because
he had already done so by the previous day's courier, ``in the
most affectionate terms''; he defends himself against Sestius's
complaint about being missed at Puteoli (Sestius was the one who
failed to come back from Cosa in time). Section 2 is a small,
moving turn: Atticus had wept on parting from him at Tusculum, and
Cicero would have changed his whole plan had he seen it. The
hope of meeting again soon is what sustains them both. He
promises letters in abundance, the book \textit{On Glory} ---
imminent --- and a smaller piece in the Heraclidean manner to lie
hidden in Atticus's library.}

\textit{Section 3 turns to the household. Cicero's grand-niece
Attica has reason to complain (of his silence, presumably);
Atticus's news about Bacchis and the garlands on the statues is
welcome, and Cicero asks that nothing be left out. The closing
line is one of those characteristic Ciceronian flashes: ``What a
disgraceful son your sister has!'' --- of Quintus the younger,
arriving uninvited at dinner --- with the Homeric phrase
\textit{aut\=ei boulysei} (``with the ox-loosing,'' i.e. at the
hour when ploughmen unyoke their oxen) to fix the moment of his
arrival in the high register of epic.}
